\section{Experimental Setup}
The source code and gold standard are publicly available at
\url{https://github.com/Mango-Cats/tagabaybay}. The rules are implemented in
Rust; the G2P backend uses the phonemizer library~\cite{bernard2021}, which
wraps eSpeak-NG.

\subsection{Gold Standard}

The gold standard corpus comprises 2,319 English--Filipino loanword pairs,
drawn from a list English words commonly used in Filipino everyday speech
(1,858 pairs) and pharmaceutical terminology (461 pairs). Three native Filipino
mesolectal speakers, familiar with the KWF manual beforehand, independently
annotated the full word list to produce nativized forms. Annotators applied the
spelling conventions prescribed in~\cite{almario2014,schachter1972}; where the
manual did not prescribe a specific nativization, they
consulted~\cite{llamzon1966,llamzon1997} to determine the most plausible form
under the five-vowel system.

Before adjudication, majority-match agreement was 92\%. The remaining items,
where no two annotators agreed, were resolved through discussion among all
three annotators, and the corpus was finalized accordingly.

\subsection{Evaluation Metrics}\label{ssec:evaluation_metrics}

The primary evaluation metric is Character Error Rate (CER), computed as the
Levenshtein edit distance between the system output and the gold standard
reference, normalized by the total character count of the reference ascross the
entire corpus:
\begin{equation}
    \text{CER} = \frac{\sum_{i} \textsc{Levenshtein}(\hat{y}_i,\, y_i)}
    {\sum_{i} |y_i|}
    \times 100\%
\end{equation}
where \(\hat{y}_i\) is the system's nativized form of the \(i\)-th word,
\(y_i\) is the gold standard reference, \(\textsc{Levenshtein}(\cdot,\cdot)\) denotes
Levenshtein distance, and \(|y_i|\) is the character length of
the reference. CER is preferred over word-level accuracy because nativization
errors are often minor, and a word-level metric would treat a one-character
discrepancy as harshly as a completely incorrect output.

Additionally, because (\gr{e}, \gr{i}, \gr{y}) and (\gr{o}, \gr{u}) are
interchangeable in Filipino~\cite{schachter1972}, we report CER under
additional conditions treating these pairs as equivalent.

\subsection{Comparison Systems}

Three additional systems are evaluated. The baseline applies a flat sequence of
regular expression substitutions without graphemic tokenization, priority
ordering, or phonetic disambiguation. The orthographic-only ablation runs the
full proposed pipeline with phonetic vowel resolution disabled; all
orthographic rules remain active, but vowel graphemes map to their identity.
The LLM prompting system sends each word individually to Claude Haiku 4.5. The
prompt provides the KWF mapping rules as a bulleted list covering consonant
correspondences, diphthong conventions, doubled-consonant simplification,
silent letter deletion, and epenthesis instructions. Each word is sent as a
single API call with no examples and no prompt engineering iteration beyond the
initial rule list.